% Generated draft: requires mathematical and source review.
\documentclass[11pt]{article}
\usepackage[margin=25mm]{geometry}
\usepackage{fontspec}
\setmainfont{Latin Modern Roman}
\usepackage{amsmath,amssymb,mathtools,unicode-math}
\setmathfont{Latin Modern Math}
\usepackage{enumitem,needspace,xurl,hyperref}
\setlength{\parindent}{0pt}
\setlength{\parskip}{5pt}
\setlength{\emergencystretch}{3em}
\newcommand{\N}{\mathbb{N}}
\newcommand{\Z}{\mathbb{Z}}
\newcommand{\Q}{\mathbb{Q}}
\newcommand{\R}{\mathbb{R}}
\newcommand{\C}{\mathbb{C}}
\newcommand{\F}{\mathbb{F}}
\newcommand{\A}{\mathbb{A}}
\newcommand{\PP}{\mathbb{P}}
\newcommand{\E}{\mathbb{E}}
\newcommand{\eps}{\varepsilon}
\newcommand{\dd}{\,\mathrm{d}}
\begin{document}
\section*{UnsolvedMath ID: 20000325}
\textbf{Source problem code:} catalogue entry
\subsection{Definitions and mathematical statement}
\textbf{Original problem (Kollár, extracted as Problem 1.6).} Consider a moduli space $M$ of canonically polarized surfaces over $\operatorname{Spec}\Z$. For a prime $p$, is its fiber $M\times_{\Z}\F_p$ the moduli space of surfaces over $\F_p$? The source says a similar result for $M_g$, the moduli of genus-$g$ curves, may be due to Oort, and asks the analogous question for other problems, including $A_g$, the moduli of principally polarized abelian varieties in the usual convention, and K3 surfaces. Polarizations, fixed numerical invariants, smooth versus stable objects, levels, and the meaning of ``moduli space'' are not specified.

For a genuinely relative, base-change-stable moduli stack $\mathcal M\to\operatorname{Spec}\Z$, its fiber $\mathcal M_p$ describes the same moduli functor restricted to $\F_p$-schemes by definition. This does not settle compatibility of coarse spaces. If $\mathcal M\to M$ is a coarse moduli morphism and $M_p^{\mathrm{coarse}}$ is the coarse space of $\mathcal M_p$, there is a comparison
\[M_p^{\mathrm{coarse}}\longrightarrow M\times_{\Z}\F_p,\]
and the substantive coarse-space question asks when it is an isomorphism. An integral model defined as a closure of characteristic-zero objects would also require comparison with the full characteristic-$p$ functor.

\textbf{Attached unreviewed research-aid claims.} For a fixed quasi-compact algebraic stack of finite presentation over $\operatorname{Spec}\Z$ with finite inertia, the aid claims uniform boundedness of stabilizer lengths, tameness after inverting finitely many primes, and compatibility of its coarse space with arbitrary base change there. Tame means that the geometric stabilizers are linearly reductive. It concludes that the displayed comparison is an isomorphism for all sufficiently large $p$ for this fixed-data stack. Applicability of these hypotheses and existence of the required coarse spaces for each named moduli problem must still be checked; no bound uniform over all numerical invariants is supplied.

The aid describes a wild-prime obstruction on finite constant-group quotient charts over a DVR. To state the reported Tor mechanism precisely in its natural flat form, let $V$ be a DVR with residue field $\kappa$, let $A$ be a flat $V$-algebra with an action of a finite constant group $G$, and put
\[d:A\longrightarrow\bigoplus_{g\in G}A,\qquad d(a)=(g(a)-a)_{g\in G},\qquad C=\operatorname{coker}d.\]
The claimed base-change defect is described by
\[0\longrightarrow A^G\otimes_V\kappa\longrightarrow(A\otimes_V\kappa)^G\longrightarrow\operatorname{Tor}_1^V(C,\kappa)\longrightarrow0.\]
The summary does not specify flatness or other precise chart hypotheses; the displayed flat interpretation must not be asserted for arbitrary nonflat charts. It also reports a proper quotient stack $[\PP^1_{\Z_{(2)}}/C_2]$ whose special-fiber coarse comparison is Frobenius on $\PP^1_{\F_2}$ and hence not an isomorphism. Its action and construction are not supplied in this record.

\textbf{Fused following-section context.} The verification field identifies the rest of the extracted text as the beginning of Section 1.4, ``Other moduli problems''. It attributes to Schumacher and Tsuji [ST04] the assertion that a separated moduli space of smooth polarized manifolds is always quasi-projective; cites Viehweg's earlier result for canonically polarized manifolds and polarized manifolds $F$ with $\omega_F$ nef; notes the difficulty of separatedness for polarized uniruled manifolds; and says Kollár proposed counterexamples by realizing every smooth toric variety as a moduli space of some polarized manifolds, using smooth proper nonprojective toric varieties. These are retained as attributed adjacent context, with no endorsement of either competing general assertion and no invented identification with the fiber question.
\subsection{Short English statement}
The problem asks whether integral moduli spaces have the expected characteristic-$p$ fibers, for surfaces and other moduli problems. Stack fibers and coarse-space base change are distinguished. The aid reports compatibility away from finitely many wild characteristics for a fixed finite-inertia stack and a quotient obstruction at two; the OCR-fused quasi-projectivity discussion is retained as adjacent context.
\subsection{Sources}
\begin{itemize}
\item Local UnsolvedMath registry, permanent ID 20000325, source code catalogue entry; \texttt{aim-algebraic-geometry-notes.json}, notes[324], tag \texttt{problem}.
\item AIM workshop \emph{Compact moduli spaces and birational geometry}; supplied verification calls the document \emph{Open Problems in Compact Moduli Spaces and Birational Geometry}, printed page 2: \url{https://aimath.org/WWN/birational/birational.pdf}. The problem is attributed to Kollár and tentatively mentions Oort. The verification describes an OCR fusion, but does not supply a full replacement transcription. No PDF was consulted here.
\item Adjacent extracted section names Schumacher and Tsuji [ST04], Viehweg, and Kollár. The bibliography for [ST04] and the other assertions is not supplied; the original literature field is empty.
\item Attached research attempt 1: result \texttt{partial\_result}, status \texttt{partially\_solved}, rigor label \texttt{proved}, novelty confidence low. It is a machine-generated, non-peer-reviewed aid requiring independent verification before citation, and supplies no theorem citations for tame base change or its quotient example.
\end{itemize}
\end{document}
